\section*{Hoofdstuk \ref{ch:b3:lab-techniques-3} --- Laboratoriumtechnieken III: onderzoekspraktijk}

\begin{solution}{exo:b3:lab-techniques-3:1}
$(0.50/1013)^3 = \num{1.2e-10}$. Lekken, gas dat op het glas geadsorbeerd is en het ontgassen van
vet en septa laten veel meer over; droog, heet glaswerk en goede
verbindingen zijn belangrijker dan extra cycli.
\end{solution}

\begin{solution}{exo:b3:lab-techniques-3:2}
$\qty{20.0}{mmol}/\qty{1.48}{mol/L} = \qty{13.5}{mL}$.
\end{solution}

\begin{solution}{exo:b3:lab-techniques-3:3}
$120 + 10 \times 1.00782503 + 55.93493633 - 0.00054858 = \qty{186.01264}{u}$. Fout: $10^6 \times (186.0129 - 186.01264)/186.01264 = +1.4$ ppm,
ruim binnen 5 ppm.
\end{solution}

\begin{solution}{exo:b3:lab-techniques-3:4}
Primair: het tijdschriftartikel, het octrooi, het proefschrift.
Secundair: het overzichtsartikel, het handboek van constanten.
Tertiair: het leerboek.
\end{solution}

\begin{solution}{exo:b3:lab-techniques-3:5}
$M = \qty{194.0}{g/mol}$: C $96.0/194.0 = 49.48$\,\%, H $10.0/194.0 = 5.15$\,\%, N $56.0/194.0 = 28.87$\,\%. Verschillen 0.17, 0.07 en 0.16
punt: binnen 0.4, de analyse ondersteunt de formule.
\end{solution}

\begin{solution}{exo:b3:lab-techniques-3:6}
$\qty{170.0}{mg}/\qty{212.0}{g/mol} = \qty{0.802}{mmol}$; $c = \qty{0.802}{mmol}/\qty{0.54}{mL} = \qty{1.5}{mol/L}$ (twee beduidende cijfers, bepaald door
het afgelezen volume).
\end{solution}

\begin{solution}{exo:b3:lab-techniques-3:7}
De gemiddelden zijn 73.0 en 77.0\,\%; beide standaardafwijkingen zijn $\sqrt{10/4} = 1.58$; $s_p = 1.58$. $t = 4.0/(1.58\sqrt{2/5}) = 4.0$, groter dan
2.31: de verbetering is significant op het niveau van 95\,\%.
\end{solution}

\begin{solution}{exo:b3:lab-techniques-3:8}
$F = (0.80/0.30)^2 = 7.1 > 5.05$: de eerste methode is significant minder precies.
\end{solution}

\begin{solution}{exo:b3:lab-techniques-3:9}
Koop kleine flessen met inhibitor, dateer elke fles bij ontvangst en bij
opening, bewaar ze gesloten, donker en koel; test op peroxiden vóór elke
destillatie of indamping en op vaste tijdstippen na opening; voer
flessen waarvan de datum verstreken is, of die positief testen, af via
de route voor gevaarlijk afval zonder ze te concentreren.
\end{solution}

\begin{solution}{exo:b3:lab-techniques-3:10}
Gevaren: vrijkomende waterstof (brand, druk), de reactiviteit van NaH met
water, en een bekende exotherme ontleding van DMF met natriumhydride die
bij verwarming op hol kan slaan; DMF is bovendien giftig voor de
voortplanting. Beheersmaatregelen: vervang het oplosmiddel (THF of
2-methyltetrahydrofuraan) of de base; als men ze behoudt, eerst op
kleine schaal, bij lage temperatuur, NaH in porties toevoegen terwijl
men de temperatuur volgt, in een zuurkast achter een scherm, met afvoer
via een olieslot en een koelbad bij de hand; een schriftelijke
procedure, een tweede persoon aanwezig, beschermingsmiddelen als
laatste.
\end{solution}

\begin{solution}{exo:b3:lab-techniques-3:11}
Relatieve bijdragen: 0.50, 0.50, $0.010/20.0 = 0.05$, $0.010/15.0 = 0.067$ en 0.10\,\%; de
gecombineerde onzekerheid is $\sqrt{0.25 + 0.25 + 0.0025 + 0.0044 + 0.010} = 0.72$\,\%. De twee integralen overheersen: een betere
signaal-ruisverhouding en een zorgvuldige basislijn- en fasecorrectie
zijn belangrijker dan de balans.
\end{solution}

\begin{solution}{exo:b3:lab-techniques-3:12}
$Q = (89.0 - 82.0)/(89.0 - 80.6) = 0.83 > 0.71$: de toets signaleert de waarde. Controleer vóór de verwerping het \omterm{def:b3:lab-techniques-3:notebook}{labjournaal}
op een oorzaak (een verkeerd afgelezen buret, een luchtbel, een andere
partij titrant); als er een gevonden wordt, verwerp de waarde en vermeld
waarom; zo niet, behoud ze, rapporteer de toets, of herhaal de meting.
\end{solution}

\begin{solution}{pb:b3:lab-techniques-3:1}
\textbf{1.} $\ce{FeCl2 + 2NaC5H5 -> Fe(C5H5)2 + 2NaCl}$.

\textbf{2.} Natriumcyclopentadienide wordt door zuurstof en water
vernietigd, en ijzer(II)chloride wordt in vochtige lucht geoxideerd tot
ijzer(III).

\textbf{3.} $\qty{5.00}{g}/\qty{126.8}{g/mol} = \qty{0.0394}{mol}$; $\times \qty{185.8}{g/mol} = \qty{7.33}{g}$; rendement $5.86/7.33 = 80$\,\%.

\textbf{4.} $(0.10/1000)^3 = 10^{-12}$ (in theorie).

\textbf{5.} Door het dimeer te verhitten, dat een retro-diels-alderreactie
ondergaat, en het vluchtige monomeer af te destilleren; bij
kamertemperatuur dimeriseert het monomeer opnieuw via een
diels-alderreactie, zodat het koud bewaard en meteen gebruikt wordt.

\textbf{6.} Elke portie zet waterstof en warmte vrij: langzaam toevoegen
houdt de gasstroom en de temperatuur onder controle.

\textbf{7.} \qty{186.0126}{u}; fout $+1.4$ ppm.

\textbf{8.} C $120.0/185.8 = 64.58$\,\%, H $10.0/185.8 = 5.38$\,\%; gevonden 64.41 en 5.47: verschillen $-0.17$ en $+0.09$, binnen 0.4.

\textbf{9.} De tien waterstofatomen zijn door de symmetrie gelijkwaardig,
en de ringen draaien snel.

\textbf{10.} Eén enkel signaal, omdat alle tien de koolstofatomen
gelijkwaardig zijn.

\textbf{11.} Een röntgenstructuur van een éénkristal.

\textbf{12.} Nee: als complex met 18 elektronen (\cref{ch:b3:organometallic-bonding}) is het stabiel aan de lucht, in
tegenstelling tot zijn precursoren.

\textbf{13.} Zijn signalen overlappen niet met die van ferroceen, het is
een zuivere, niet-vluchtige, niet-hygroscopische vaste stof die
nauwkeurig gewogen kan worden, en het is inert tegenover het monster.

\textbf{14.} Volledige relaxatie van elk signaal tussen de scans (een
wachttijd van meerdere $T_1$), een uniforme excitatie, voldoende
signaal-ruisverhouding, en een zorgvuldige fase- en basislijncorrectie.

\textbf{15.} Ferroceen $\ce{C10H10Fe}$ \qty{185.8}{g/mol}; 1,3,5-trimethoxybenzeen $\ce{C9H12O3}$ \qty{168.0}{g/mol}.

\textbf{16.} $P_x = 3.942 \times (3/10) \times (185.8/168.0) \times (15.0/20.0) \times 0.999 = 0.980$: 98.0\,\%.

\textbf{17.} Relatief $\sqrt{0.50^2 + 0.50^2 + 0.05^2 + 0.067^2 + 0.10^2} = 0.72$\,\%; absoluut $0.72\,\% \times 98.0 = 0.7$ procentpunt.

\textbf{18.} Ja, maar zwak: een onzuiverheid van 2\,\% met een
gelijkaardige samenstelling verandert C en H minder dan de analyse kan
detecteren; qNMR is de informatievere zuiverheidsmeting.

\textbf{19.} Natriumhydride (waterstof met water, brand), de vrijkomende
waterstof, THF en cyclopentadieen (zeer licht ontvlambaar; THF een
\omterm{def:b3:lab-techniques-3:peroxide}{peroxidevormer}), het hete kraken van het dimeer, ijzer(II)chloride
(irriterend), en ferroceen (ontvlambare vaste stof, schadelijk).

\textbf{20.} Natriumhydride: afgewogen als dispersie in olie onder inert
gas, in porties toegevoegd aan een geroerde, gekoelde oplossing, het gas
afgevoerd via een olieslot, geen water in de buurt. THF: droog en
peroxidevrij uit een kolom, gebruikt in de zuurkast onder stikstof.

\textbf{21.} Onder inert gas en met koeling, door langzaam isopropanol
toe te voegen, daarna ethanol, daarna water, en telkens te wachten tot
de gasontwikkeling stopt.

\textbf{22.} Een reversibele golf met één elektron, $\ce{Fe(C5H5)2+}/\ce{Fe(C5H5)2}$, met pieken die bij \qty{25}{\celsius} ongeveer 59 mV
uit elkaar liggen
(\cref{ch:b3:electrode-kinetics}); het koppel is snel, stabiel en bijna
ongevoelig voor het oplosmiddel, wat er een handige interne referentie
van maakt.

\textbf{23.} Het \omterm{def:b3:lab-techniques-3:notebook}{labjournaal}: datum, schema, hoeveelheden en herkomst, de
\omterm{def:b3:lab-techniques-3:assessment}{risicobeoordeling}, wat gedaan en gezien werd, rendementen en de
bestandsnamen van de \omterm{def:b3:lab-techniques-3:notebook}{ruwe gegevens}. De \omterm{def:b3:lab-techniques-3:literature}{ondersteunende informatie}: de
volledige procedure en alle karakteriseringsgegevens met kopieën van de
spectra.

\textbf{24.} Zodat anderen, en de auteurs jaren later, ze kunnen openen en
opnieuw verwerken, met welke software ze ook hebben: vindbaar,
toegankelijk, interoperabel en herbruikbaar.

\textbf{25.} De qNMR-zuiverheid van het ferroceenmonster is $98.0 \pm 0.7$\,\% (standaardonzekerheid).
\end{solution}
